\section{Introduction}
\label{sec:intro}

Visual Question Answering (VQA) asks a model to answer a natural-language question
about an image, which requires it to perceive the visual content, understand the
question and reason over both. The task has progressed rapidly with large
vision--language models (VLMs) such as BLIP-2~\cite{li2023blip2},
Flamingo~\cite{alayrac2022flamingo} and LLaVA~\cite{liu2023visualinstructiontuning},
which connect a pre-trained vision encoder to a pre-trained language model and
generate free-form answers. Vietnamese has followed this trend: Vintern-1B~\cite{doan2024vintern1b}
connects an InternViT-300M encoder~\cite{chen2024internvl} to a Qwen2.5-0.5B
decoder~\cite{qwen2025qwen25} and is among the strongest open Vietnamese VLMs, while
benchmarks such as ViVQA~\cite{tran2021vivqa}, OpenViVQA~\cite{nguyen2023openvivqa} and
the recent reasoning-oriented AutoViVQA~\cite{autovivqa2026} make it possible to measure
progress in this low-resource setting.

Obtaining a strong model on a new Vietnamese benchmark, however, is still expensive.
The common recipe fine-tunes an existing VLM: the fine-tuned Vintern-1B baseline of
AutoViVQA adds LoRA adapters to the language model and splits every image into up to
six high-resolution tiles plus a thumbnail, so that the decoder reads up to 1{,}792
visual tokens per image. The alternative is to design a task-specific architecture, as
ViMoE-VQA~\cite{vimoevqa} does with a sparse mixture-of-experts fusion module and a
decoder trained on AutoViVQA. A cheaper option is to keep the whole VLM frozen and train
only a small connector between its encoder and decoder, but the design of that
connector is delicate. A connector that reads only the global image vector is light,
yet it cannot convey small objects, printed text or spatial relations. Connectors that
learn to compress the patch tokens, such as Q-Formers, can convey such details but
must learn a new alignment to the decoder from limited Vietnamese data.

Vintern-1B already contains a component that solves this alignment problem: its
projector, which was pre-trained to map image regions into the input space of the
decoder. At the same time, it discards the global \texttt{[CLS]} vector of its encoder,
which summarises the whole scene. This observation raises a simple question:
\emph{which visual tokens should a frozen Vietnamese VLM receive, and how many?}

To address this question, we propose \textbf{ViBridge-VQA}, which gives the frozen
decoder of Vintern-1B two complementary kinds of visual tokens from a single image view.
\emph{Local} tokens reuse the frozen projector of Vintern-1B, average-pooled to a chosen
budget and placed in its original image slot, so the decoder receives region features
in the representation it was trained to read. \emph{Global} tokens are produced from the
\texttt{[CLS]} vector by a small trainable bridge and add the scene-level information that
Vintern-1B ignores. Only this bridge is trained. Because every additional token
increases inference cost, we treat the number of global and local tokens as a budget
and select it as the knee of the accuracy--cost Pareto front. As a result, ViBridge-VQA
with 14 global and 36 local tokens outperforms the officially fine-tuned Vintern-1B on most metrics while reading up to 36 times fewer visual tokens and updating none of its weights.

In summary, our main contributions are as follows:
\begin{itemize}
  \item We propose \textbf{ViBridge-VQA}, a global--local bridge that improves a frozen
  Vietnamese VLM by combining trainable global tokens with local tokens from its own
  pre-trained projector. It raises the F1 of Vintern-1B on AutoViVQA from 17.6
  (zero-shot) to 54.6 with 1.35\% of the parameters trained, and to 58.2 when combined
  with a lightweight decoder LoRA.
  \item We study the \textbf{visual token budget} systematically over eight
  configurations and three random seeds, and select it with an explicit criterion. The
  knee of the Pareto front falls on the same configuration for every metric, on both
  evaluation splits, and for cost measured in tokens or in inference time.
  \item We \textbf{analyse where the gains come from}: local tokens matter far more than
  global ones, should come from the pre-trained projector, complement decoder
  adaptation, and transfer to four unseen Vietnamese VQA datasets.
\end{itemize}
